% BHU PYQ -- Manually transcribed & error-corrected
% Paper: BCH-215 : Fundamentals of Business Finance
% Exam: B.Com. (Hons.) Semester III Examination, 2022-23

\documentclass[11pt,a4paper]{article}
\usepackage[a4paper,top=2.5cm,bottom=2.5cm,left=2.8cm,right=2.8cm,headheight=14pt]{geometry}
\usepackage{amsmath,amssymb}
\usepackage[T1]{fontenc}
\usepackage[utf8]{inputenc}
\usepackage{lmodern,microtype,fancyhdr,enumitem,hyperref,lastpage,parskip,booktabs}

\hypersetup{
    colorlinks=true,
    urlcolor=black,
    linkcolor=black,
    pdftitle={BCH-215: Fundamentals of Business Finance -- BHU B.Com. (Hons.) Sem III 2022-23}
}

\pagestyle{fancy}
\fancyhf{}
\renewcommand{\headrulewidth}{0.4pt}
\renewcommand{\footrulewidth}{0.4pt}
\fancyhead[L]{\small   \textit{Banaras Hindu University}}
\fancyhead[R]{\small   \textit{BCH-215: Fundamentals of Business Finance}}
\fancyfoot[C]{\small Page  \thepage\ of \pageref{LastPage}}


\newlist{parts}{enumerate}{2}
\setlist[parts,1]{label=(\alph*),leftmargin=*,itemsep=5pt,topsep=3pt}
\setlist[parts,2]{label=(\roman*),leftmargin=*,itemsep=3pt,topsep=2pt}

\newcommand{\pts}[1]{\hfill   \textit{[#1]}}


\begin{document}


\begin{center}
  {\large\bfseries BANARAS HINDU UNIVERSITY}\\[2pt]
  {\normalsize B.Com. (Hons.) --- Semester III Examination, 2022-23}\\[6pt]
  \rule{0.8\linewidth}{0.5pt}\\[6pt]
  {\large\bfseries Paper No. BCH-215: Fundamentals of Business Finance}\\[4pt]
  {\normalsize Time: 3 Hours\hfill Maximum Marks: 70}
\end{center}

\rule{\linewidth}{0.4pt}

\smallskip



\noindent   \textit{Instructions: Answer all questions. All questions carry equal marks (14 marks each).}

\bigskip




\noindent   \textbf{Question 1.} \\
Point out the meaning and scope of finance function. ``The financial man should not be only money man but also a business man.'' Elucidate the statement.\pts{14}

\centerline{   \textbf{OR}}

How is wealth maximisation ideal a conceptual refinement of the profit maximisation idea? How can the former be meaningfully termed in a volatile financial market? Explain.\pts{14}

\medskip\rule{\linewidth}{0.2pt}




\noindent   \textbf{Question 2.} \\
Mention the concept and kinds of financial planning. What are the basic characteristics of a sound financial plan for a joint-stock company?\pts{14}

\centerline{   \textbf{OR}}

``You are requested to conduct the financial planning for a new manufacturing company to be set up. What factors would you keep in mind while doing so?''\pts{14}

\medskip\rule{\linewidth}{0.2pt}




\noindent   \textbf{Question 3.} \\
What is capitalisation? Briefly explain its theories.\pts{14}

\centerline{   \textbf{OR}}


\begin{parts}
  \item What are the causes of watered capital? How does it differ from overcapitalisation?\pts{7}
  \item Three companies whose Balance Sheets are summarised below are engaged in a similar business. If their current rate of capitalisation is 8\%, analyse the state of capitalisation of each company:
  
  
\begin{center}
    \small
    
\begin{tabular}{lrrr}
       oprule
   \textbf{Particulars} &   \textbf{Company A (Rs.)} &   \textbf{Company B (Rs.)} &   \textbf{Company C (Rs.)} \\
      \midrule
      Share capital @ Rs. 10 & 3,00,000 & 2,00,000 & 4,00,000 \\
      Reserve and surplus & 1,00,000 & 1,00,000 & 1,00,000 \\
      Current liabilities & 50,000 & 60,000 & 80,000 \\
      Sundry assets & 4,50,000 & 3,60,000 & 5,80,000 \\
      Average profits & 30,000 & 24,000 & 42,000 \\
      ottomrule
    \end{tabular}
  \end{center}\pts{7}
\end{parts}

\medskip\rule{\linewidth}{0.2pt}




\noindent   \textbf{Question 4.} \\
What is the meaning of share? How is it different from stock? Briefly state the various types of shares.\pts{14}

\centerline{   \textbf{OR}}

Critically examine the role of financing through debentures. Why has it not been popular in India?\pts{14}

\medskip\rule{\linewidth}{0.2pt}




\noindent   \textbf{Question 5.} \\
Define working capital and discuss its determining factors.\pts{14}

\centerline{   \textbf{OR}}

Write lucid notes on any two of the following:\pts{7+7}

\begin{parts}
  \item Bonds and term loans
  \item Undercapitalisation
  \item Trading on equity
\end{parts}

\vfill
\rule{\linewidth}{0.4pt}

\begin{center}\small
  \href{https://bhu-exam.vercel.app/}{Click here to visit our website}\\[4pt]
  \href{https://me-aryan.vercel.app/}{Click here to visit our contributor}\\[4pt]
  \href{https://quashan-me.vercel.app/}{Click here to visit our contributor}
\end{center}

\end{document}
